% Fuente autónoma generada desde los contenidos compartidos del informe.

% Compilar con pdfLaTeX; no necesita archivos de imagen ni shell-escape.

\documentclass[10pt,a4paper]{article}
\usepackage[T1]{fontenc}
\usepackage[utf8]{inputenc}
\usepackage{lmodern}
\usepackage[spanish,es-nodecimaldot,es-noshorthands]{babel}
\usepackage[a4paper,left=14mm,right=14mm,top=14mm,bottom=16mm,footskip=9mm]{geometry}
\usepackage[table]{xcolor}
\usepackage{array,tabularx}
\usepackage{amsmath,amssymb}
\usepackage{tikz}
\usepackage{fancyhdr}
\usepackage{enumitem}
\usepackage{url}
\usepackage[unicode=true,pdfencoding=auto,colorlinks=true,urlcolor=Accent,linkcolor=Accent]{hyperref}
\definecolor{Accent}{HTML}{194B7B}
\definecolor{BodyColor}{HTML}{213046}
\definecolor{Muted}{HTML}{52677E}
\definecolor{CodeBg}{HTML}{F0F4F9}
\definecolor{TableHead}{HTML}{E9F0F8}
\definecolor{RuleColor}{HTML}{CAD5E1}
\definecolor{NodeLine}{HTML}{51749B}
\definecolor{NodeBg}{HTML}{EDF3FB}
\definecolor{LeafBg}{HTML}{E9F5ED}
\definecolor{EdgeColor}{HTML}{5A6C80}
\renewcommand{\familydefault}{\sfdefault}
\urlstyle{same}
\setlength{\parindent}{0pt}
\setlength{\parskip}{5pt}
\setlength{\emergencystretch}{2em}
\clubpenalty=10000
\widowpenalty=10000
\raggedbottom
\pagestyle{fancy}
\fancyhf{}
\renewcommand{\headrulewidth}{0pt}
\renewcommand{\footrulewidth}{.35pt}
\fancyfoot[L]{\fontsize{8}{10}\selectfont\color{Muted}INFO1148 \textperiodcentered{} Laboratorio 3 \textperiodcentered{} Propuesta de solución}
\fancyfoot[R]{\fontsize{8}{10}\selectfont\color{Muted}\thepage}
\newcommand{\ReportSection}[1]{\par{\color{Accent}\bfseries\fontsize{16}{20}\selectfont #1\par}\vspace{5pt}}
\newcommand{\ReportHeading}[1]{\par\addvspace{6pt}{\color{Accent}\bfseries\fontsize{10.8}{14}\selectfont #1\par}\nobreak\vspace{2pt}}


\hypersetup{pdftitle={Laboratorio 3 - Diseño y Validación de Gramáticas},pdfauthor={Juan Cárdenas y Bayron Flores}}

\begin{document}

\color{BodyColor}\fontsize{9.2}{13.0}\selectfont

\begin{center}
{\color{black}\bfseries\fontsize{17}{21}\selectfont INFORME DE LABORATORIO \#3\par}
\vspace{4pt}
{\color{black}\bfseries\fontsize{15}{19}\selectfont Diseño y Validación de Gramáticas\par}
\end{center}
\vspace{9pt}

\begingroup\fontsize{10}{13}\selectfont
\setlength{\tabcolsep}{7pt}\renewcommand{\arraystretch}{1.3}
\arrayrulecolor{black}
\noindent\begin{tabularx}{\linewidth}{|>{\raggedright\arraybackslash}p{0.38\dimexpr\linewidth-4\tabcolsep-3\arrayrulewidth\relax}|>{\raggedright\arraybackslash}X|}\hline
\rowcolor{TableHead}\centering\textbf{Antecedente} & \centering\arraybackslash\textbf{Información} \\\hline
Asignatura & Teoría de la Computación INFO1148 \\\hline
Carrera & Ingeniería Civil en Informática \\\hline
Docente & M. Lévano \\\hline
Integrantes & Juan Cárdenas y Bayron Flores (Jefe de Grupo) \\\hline
Fecha & 21 de octubre de 2026 \\\hline
\end{tabularx}\par\endgroup
\vspace{10pt}

\ReportSection{1. Introducción}

El propósito de este laboratorio es diseñar y validar una Gramática Libre de Contexto (GLC) para un lenguaje de programación simplificado. Se identifican sus terminales, no terminales y símbolo inicial; se construyen producciones para programas con una o más instrucciones, asignaciones y salidas; y se comprueba su coherencia mediante derivaciones por la izquierda, un árbol sintáctico y cadenas válidas e inválidas. El análisis léxico reconoce identificadores y números, mientras que el análisis sintáctico verifica el orden de sus categorías. Como extensión creativa se incorporan declaraciones de variables de tipo entero y real. El informe presenta las decisiones de diseño y las evidencias que permiten verificar cada resultado.
\par

\ReportSection{2. Desarrollo}
\ReportHeading{2.1 Pregunta 1 - Identificación formal de la gramática}

\ReportHeading{1. Terminales y clasificación}

Separamos el análisis léxico del sintáctico. El léxico reconoce nombres y números concretos; la GLC trabaja con sus categorías ID y NUM. Los otros terminales se escriben por su lexema literal.
\par

\begingroup\fontsize{8.0}{10.8}\selectfont
\setlength{\tabcolsep}{4pt}\renewcommand{\arraystretch}{1.16}
\arrayrulecolor{RuleColor}
\noindent\begin{tabularx}{\linewidth}{|>{\hsize=0.600000\hsize\linewidth=\hsize\raggedright\arraybackslash}X|>{\hsize=0.810000\hsize\linewidth=\hsize\raggedright\arraybackslash}X|>{\hsize=1.590000\hsize\linewidth=\hsize\raggedright\arraybackslash}X|}\hline
\rowcolor{TableHead} \textbf{Clase} & \textbf{Terminales} & \textbf{Función} \\\hline
Palabras reservadas & Begin, Programa, End, salida & Delimitan el programa o introducen una salida. \\\hline
Operadores aritméticos & +, -, *, /\allowbreak{}, \% & Operaciones binarias entre dos operandos. \\\hline
Asignación & = & Separa el identificador destino de su valor o expresión. \\\hline
Símbolo especial & \# & Termina una instrucción. \\\hline
Identificadores & ID & Categoría de nombres como x, z2 y operativo\_\allowbreak{}01\_\allowbreak{}02. \\\hline
Números & NUM & Categoría de valores como 20, 0 y 1.345. \\\hline
\end{tabularx}\par\endgroup\vspace{5pt}

\ReportHeading{2. No terminales y 3. símbolo inicial}

V = \{\textless{}programa\textgreater{}, \textless{}lista\_\allowbreak{}instrucciones\textgreater{}, \textless{}instruccion\textgreater{}, \textless{}asignacion\_\allowbreak{}simple\textgreater{}, \textless{}asignacion\_\allowbreak{}aritmetica\textgreater{}, \textless{}expresion\textgreater{}, \textless{}operando\textgreater{}, \textless{}operador\textgreater{}, \textless{}salida\textgreater{}\}. El símbolo inicial es S = \textless{}programa\textgreater{}.
\par

\ReportHeading{4. Definición formal}

\begingroup\setlength{\fboxsep}{6pt}
\noindent\colorbox{CodeBg}{%
\begin{minipage}{\dimexpr\linewidth-2\fboxsep\relax}
\ttfamily\fontsize{8.0}{10.5}\selectfont\raggedright
\setlength{\parindent}{0pt}\setlength{\parskip}{0pt}
\strut{}G = (V, \ensuremath{\Sigma}, P, S)\par
\strut{}\ensuremath{\Sigma} = \{Begin, Programa, End, salida, ID, NUM, =, \#, +, -, *, /, \%\}\par
\strut{}S = \textless{}programa\textgreater{}\par
\end{minipage}}\par\endgroup\vspace{5pt}

V contiene los nueve no terminales anteriores y \ensuremath{\Sigma} los trece terminales indicados. V y \ensuremath{\Sigma} son disjuntos. P es el conjunto de producciones de la sección 2. Es una GLC porque el lado izquierdo de cada producción es un único no terminal.
\par

\textbf{Notación:} se conserva el orden G = (V, \ensuremath{\Sigma}, P, S) pedido por el laboratorio. En otros apuntes del curso la misma información se ordena como (V\_\allowbreak{}N, V\_\allowbreak{}T, S, P). Aquí V significa el conjunto de no terminales.
\par

\textit{Respaldo: enunciado, PDF p. 2, preguntas 1-5; sesion\_\allowbreak{}3\_\allowbreak{}4.pdf, PDF pp. 3 y 8.}
\par

\clearpage

\ReportSection{Identificadores, números y decisiones léxicas}

\ReportHeading{5. Por qué no enumerar cada nombre o valor}

Sería insuficiente listar únicamente x, y, z o 20: también se necesitan nombres nuevos y números de longitud arbitraria. Se usan patrones generales, y cada token conserva el lexema concreto como dato asociado. ID y NUM son terminales de categoría de esta GLC; no se añaden reglas ID \ensuremath{\to} x ni NUM \ensuremath{\to} 20.
\par

\begingroup\setlength{\fboxsep}{6pt}
\noindent\colorbox{CodeBg}{%
\begin{minipage}{\dimexpr\linewidth-2\fboxsep\relax}
\ttfamily\fontsize{8.0}{10.5}\selectfont\raggedright
\setlength{\parindent}{0pt}\setlength{\parskip}{0pt}
\strut{}L = [A-Za-zÑñÁÉÍÓÚÜáéíóúü]\par
\strut{}D = [0-9]\par
\strut{}ID = L (L \textbar{} D \textbar{} "\_")*\par
\strut{}NUM = D+ ("." D+)?\par
\end{minipage}}\par\endgroup\vspace{5pt}

\begin{itemize}[leftmargin=13pt,itemsep=3pt,topsep=3pt,parsep=0pt]

\item Un identificador empieza con una letra; después admite letras, dígitos y guion bajo. Ejemplos: x, z2 y operativo\_\allowbreak{}01\_\allowbreak{}02. La clase L declarada es la convención finita de letras de esta implementación.

\item Las reservadas exactas son Begin, Programa, End y salida. Se distinguen mayúsculas y minúsculas. Un nombre más largo, como salida\_\allowbreak{}total, sigue siendo un identificador.

\item NUM admite enteros y reales decimales sin signo: 20, 0 y 1.345. Aunque el enunciado usa la palabra positivos, el programa incluye 0; por ello se interpretan como valores no negativos.

\item Se requieren dígitos a ambos lados del punto de un real. No se incluyen notación exponencial ni signos unarios.

\item Los espacios separan componentes y se descartan. Se admiten comentarios /\allowbreak{}* ... */\allowbreak{} no anidados para representar el comentario del programa de referencia. El comentario no se convierte en una instrucción.

\item En modo de extensión, entero y real se añaden al conjunto de reservadas. real es una palabra de tipo; NUM sigue siendo la categoría de un literal numérico.

\end{itemize}

\ReportHeading{Una categoría, muchos lexemas}

\begingroup\fontsize{8.0}{10.8}\selectfont
\setlength{\tabcolsep}{4pt}\renewcommand{\arraystretch}{1.16}
\arrayrulecolor{RuleColor}
\noindent\begin{tabularx}{\linewidth}{|>{\hsize=0.800000\hsize\linewidth=\hsize\raggedright\arraybackslash}X|>{\hsize=1.200000\hsize\linewidth=\hsize\raggedright\arraybackslash}X|}\hline
\rowcolor{TableHead} \textbf{Texto} & \textbf{Lectura sintáctica} \\\hline
x = 20 \# & ID[x] = NUM[20] \# \\\hline
z = x + y \# & ID[z] = ID[x] + ID[y] \# \\\hline
salida temporal \# & salida ID[temporal] \# \\\hline
\end{tabularx}\par\endgroup\vspace{5pt}

En ID[z], el símbolo terminal es ID y z es su lexema anotado. Esta notación no añade un nuevo terminal por cada nombre. Los corchetes son explicativos y no se escriben en el programa.
\par

\textit{Respaldo: enunciado, PDF pp. 1-2; los patrones y la separación de fases son decisiones explícitas de esta propuesta.}
\par

\clearpage

\ReportSection{2.2 Pregunta 2 - Construcción de la GLC}

\ReportHeading{6. Producciones de la gramática base}

\begingroup\setlength{\fboxsep}{6pt}
\noindent\colorbox{CodeBg}{%
\begin{minipage}{\dimexpr\linewidth-2\fboxsep\relax}
\ttfamily\fontsize{8.0}{10.5}\selectfont\raggedright
\setlength{\parindent}{0pt}\setlength{\parskip}{0pt}
\strut{}\textless{}programa\textgreater{} ::= "Begin" "Programa" ID \textless{}lista\_instrucciones\textgreater{} "End"\par
\strut{}\textless{}lista\_instrucciones\textgreater{} ::= \textless{}instruccion\textgreater{} \textless{}lista\_instrucciones\textgreater{} \textbar{} \textless{}instruccion\textgreater{}\par
\strut{}\textless{}instruccion\textgreater{} ::= \textless{}asignacion\_simple\textgreater{} \textbar{} \textless{}asignacion\_aritmetica\textgreater{} \textbar{} \textless{}salida\textgreater{}\par
\strut{}\textless{}asignacion\_simple\textgreater{} ::= ID "=" \textless{}operando\textgreater{} "\#"\par
\strut{}\textless{}asignacion\_aritmetica\textgreater{} ::= ID "=" \textless{}expresion\textgreater{} "\#"\par
\strut{}\textless{}expresion\textgreater{} ::= \textless{}operando\textgreater{} \textless{}operador\textgreater{} \textless{}operando\textgreater{}\par
\strut{}\textless{}operando\textgreater{} ::= ID \textbar{} NUM\par
\strut{}\textless{}operador\textgreater{} ::= "+" \textbar{} "-" \textbar{} "*" \textbar{} "/" \textbar{} "\%"\par
\strut{}\textless{}salida\textgreater{} ::= "salida" ID "\#"\par
\end{minipage}}\par\endgroup\vspace{5pt}

\ReportHeading{7. Diferencia entre las tres instrucciones}

\begingroup\fontsize{8.0}{10.8}\selectfont
\setlength{\tabcolsep}{4pt}\renewcommand{\arraystretch}{1.16}
\arrayrulecolor{RuleColor}
\noindent\begin{tabularx}{\linewidth}{|>{\hsize=0.720000\hsize\linewidth=\hsize\raggedright\arraybackslash}X|>{\hsize=1.440000\hsize\linewidth=\hsize\raggedright\arraybackslash}X|>{\hsize=0.840000\hsize\linewidth=\hsize\raggedright\arraybackslash}X|}\hline
\rowcolor{TableHead} \textbf{Construcción} & \textbf{Forma} & \textbf{Ejemplo} \\\hline
Asignación simple & ID = \textless{}operando\textgreater{} \# & temporal = x \# \\\hline
Asignación aritmética & ID = \textless{}operando\textgreater{} \textless{}operador\textgreater{} \textless{}operando\textgreater{} \# & z = x + y \# \\\hline
Salida & salida ID \# & salida z \# \\\hline
\end{tabularx}\par\endgroup\vspace{5pt}

Las alternativas son distinguibles: después del primer operando, \# cierra una asignación simple; un operador seguido de otro operando forma una asignación aritmética. salida inicia la tercera clase.
\par

\ReportHeading{8. Función de los no terminales}

\begingroup\fontsize{8.0}{10.8}\selectfont
\setlength{\tabcolsep}{4pt}\renewcommand{\arraystretch}{1.16}
\arrayrulecolor{RuleColor}
\noindent\begin{tabularx}{\linewidth}{|>{\hsize=0.680000\hsize\linewidth=\hsize\raggedright\arraybackslash}X|>{\hsize=1.320000\hsize\linewidth=\hsize\raggedright\arraybackslash}X|}\hline
\rowcolor{TableHead} \textbf{No terminal} & \textbf{Función} \\\hline
\textless{}programa\textgreater{} & Construye un programa completo con cabecera, cuerpo y cierre. \\\hline
\textless{}lista\_\allowbreak{}instrucciones\textgreater{} & Produce una o más instrucciones, conservando su orden. \\\hline
\textless{}instruccion\textgreater{} & Selecciona una asignación simple, una aritmética o una salida. \\\hline
\textless{}asignacion\_\allowbreak{}simple\textgreater{} & Asigna un identificador o un número a un identificador destino. \\\hline
\textless{}asignacion\_\allowbreak{}aritmetica\textgreater{} & Asigna el resultado sintáctico de una operación binaria. \\\hline
\textless{}expresion\textgreater{} & Une un operando izquierdo, un operador y un operando derecho. \\\hline
\textless{}operando\textgreater{} & Selecciona ID o NUM. \\\hline
\textless{}operador\textgreater{} & Selecciona uno de los cinco operadores aritméticos. \\\hline
\textless{}salida\textgreater{} & Forma una instrucción salida ID \#. \\\hline
\end{tabularx}\par\endgroup\vspace{5pt}

La lista nunca produce \ensuremath{\varepsilon}: tiene la alternativa de una instrucción y otra que agrega una instrucción delante de una lista. Así permite cualquier cantidad positiva de instrucciones.
\par

\textit{Respaldo: enunciado, PDF pp. 2-3, pregunta 2.}
\par

\clearpage

\ReportSection{Verificación del programa de referencia}

\ReportHeading{9. El programa completo queda representado}

La cabecera se reconoce como Begin Programa ID[operativo\_\allowbreak{}01\_\allowbreak{}02], la lista contiene veinte instrucciones y el cierre es End. Tras descartar el comentario y los espacios, cada instrucción coincide con una de las tres alternativas de la gramática.
\par

\begingroup\fontsize{8.0}{10.8}\selectfont
\setlength{\tabcolsep}{4pt}\renewcommand{\arraystretch}{1.16}
\arrayrulecolor{RuleColor}
\noindent\begin{tabularx}{\linewidth}{|>{\hsize=0.810000\hsize\linewidth=\hsize\raggedright\arraybackslash}X|>{\hsize=0.390000\hsize\linewidth=\hsize\raggedright\arraybackslash}X|>{\hsize=1.800000\hsize\linewidth=\hsize\raggedright\arraybackslash}X|}\hline
\rowcolor{TableHead} \textbf{Clase} & \textbf{Cantidad} & \textbf{Evidencia del programa} \\\hline
Asignaciones simples & 6 & x = 1.345 \#; y = 20 \#; temporal = 0 \#; temporal = x \#; x = y \#; y = temporal \# \\\hline
Asignaciones aritméticas & 7 & z y z2 usan +; h y h2 usan -; k usa *; t usa /\allowbreak{}; m usa \%. \\\hline
Salidas & 7 & Se imprimen z, h, k, t, m, x e y. \\\hline
Total & 20 & Las veinte instrucciones terminan en \#. \\\hline
\end{tabularx}\par\endgroup\vspace{5pt}

El validador acepta el archivo completo ejemplos/\allowbreak{}programa\_\allowbreak{}referencia.txt desde el símbolo inicial \textless{}programa\textgreater{}. La evidencia registra 91 tokens y los veinte nodos \textless{}instruccion\textgreater{} del árbol.
\par

\ReportHeading{Alcance elegido}

La asignación aritmética contiene exactamente una operación binaria. Esta elección sencilla cubre las estructuras mínimas y todos los ejemplos pedidos. No se agregan paréntesis ni operaciones encadenadas; por ello no hace falta fijar precedencia entre varios operadores de una misma expresión.
\par

\begingroup\fontsize{8.0}{10.8}\selectfont
\setlength{\tabcolsep}{4pt}\renewcommand{\arraystretch}{1.16}
\arrayrulecolor{RuleColor}
\noindent\begin{tabularx}{\linewidth}{|>{\hsize=1.290000\hsize\linewidth=\hsize\raggedright\arraybackslash}X|>{\hsize=0.510000\hsize\linewidth=\hsize\raggedright\arraybackslash}X|>{\hsize=1.200000\hsize\linewidth=\hsize\raggedright\arraybackslash}X|}\hline
\rowcolor{TableHead} \textbf{Cadena adicional} & \textbf{Resultado} & \textbf{Razón} \\\hline
x = a + b * c \# & Inválida & Contiene dos operadores en una sola expresión. \\\hline
x = -1 \# & Inválida & No se definió signo unario. \\\hline
Begin Programa prueba End & Inválida & La lista debe contener al menos una instrucción. \\\hline
\end{tabularx}\par\endgroup\vspace{5pt}

Este es un análisis sintáctico. No se calcula el valor de k /\allowbreak{} z ni se comprueba si una variable fue declarada o si una división puede dar error al ejecutarse.
\par

\textit{Respaldo: programa de referencia del enunciado, PDF pp. 1-2; evidencia ejecutada en evidencias/\allowbreak{}validacion.json.}
\par

\clearpage

\ReportSection{2.3 Pregunta 3 - Derivaciones y árbol sintáctico}

En cada paso se sustituye el no terminal situado más a la izquierda. Como S = \textless{}programa\textgreater{} genera programas completos, las instrucciones aisladas de los puntos 10 y 11 se derivan localmente desde \textless{}instruccion\textgreater{}.
\par

\ReportHeading{10. Derivación de z = x + y \#}

\begingroup\setlength{\fboxsep}{6pt}
\noindent\colorbox{CodeBg}{%
\begin{minipage}{\dimexpr\linewidth-2\fboxsep\relax}
\ttfamily\fontsize{8.0}{10.5}\selectfont\raggedright
\setlength{\parindent}{0pt}\setlength{\parskip}{0pt}
\strut{}\textless{}instruccion\textgreater{}\par
\strut{}\ensuremath{\Rightarrow} \textless{}asignacion\_aritmetica\textgreater{}\par
\strut{}\ensuremath{\Rightarrow} ID[z] = \textless{}expresion\textgreater{} \#\par
\strut{}\ensuremath{\Rightarrow} ID[z] = \textless{}operando\textgreater{} \textless{}operador\textgreater{} \textless{}operando\textgreater{} \#\par
\strut{}\ensuremath{\Rightarrow} ID[z] = ID[x] \textless{}operador\textgreater{} \textless{}operando\textgreater{} \#\par
\strut{}\ensuremath{\Rightarrow} ID[z] = ID[x] + \textless{}operando\textgreater{} \#\par
\strut{}\ensuremath{\Rightarrow} ID[z] = ID[x] + ID[y] \#\par
\end{minipage}}\par\endgroup\vspace{5pt}

\textbf{Realización léxica de la última fila:} z = x + y \#. Las elecciones z, x e y son los valores de los tokens ID; no son nuevas producciones.
\par

\ReportHeading{11. Derivación de salida z \#}

\begingroup\setlength{\fboxsep}{6pt}
\noindent\colorbox{CodeBg}{%
\begin{minipage}{\dimexpr\linewidth-2\fboxsep\relax}
\ttfamily\fontsize{8.0}{10.5}\selectfont\raggedright
\setlength{\parindent}{0pt}\setlength{\parskip}{0pt}
\strut{}\textless{}instruccion\textgreater{}\par
\strut{}\ensuremath{\Rightarrow} \textless{}salida\textgreater{}\par
\strut{}\ensuremath{\Rightarrow} salida ID[z] \#\par
\end{minipage}}\par\endgroup\vspace{5pt}

\textbf{Realización léxica:} salida z \#.
\par

\ReportHeading{Cómo se integran en una derivación desde S}

\begingroup\setlength{\fboxsep}{6pt}
\noindent\colorbox{CodeBg}{%
\begin{minipage}{\dimexpr\linewidth-2\fboxsep\relax}
\ttfamily\fontsize{8.0}{10.5}\selectfont\raggedright
\setlength{\parindent}{0pt}\setlength{\parskip}{0pt}
\strut{}\textless{}programa\textgreater{}\par
\strut{}\ensuremath{\Rightarrow} Begin Programa ID[prueba] \textless{}lista\_instrucciones\textgreater{} End\par
\strut{}\ensuremath{\Rightarrow} Begin Programa ID[prueba] \textless{}instruccion\textgreater{} End\par
\end{minipage}}\par\endgroup\vspace{5pt}

A partir de esa última forma se aplica cualquiera de las derivaciones locales anteriores. El resultado es un programa completo de una instrucción. Una instrucción aislada es un fragmento válido; no se afirma que por sí sola pertenezca al lenguaje completo generado desde \textless{}programa\textgreater{}.
\par

\textit{Respaldo: enunciado, PDF p. 3, puntos 10-11; las derivaciones se calculan con las producciones de la sección 2.}
\par

\clearpage

\ReportSection{Árbol sintáctico y tres cadenas adicionales}

\ReportHeading{12. Árbol de z = x + y \#}

\begin{center}
\begin{tikzpicture}[x=1cm,y=1cm,scale=1.0000,transform shape,
every node/.style={draw=NodeLine,rounded corners=2pt,inner xsep=4pt,inner ysep=4pt,font=\ttfamily\fontsize{8}{10}\selectfont},
noTerminal/.style={fill=NodeBg},terminal/.style={fill=LeafBg},
every path/.style={draw=EdgeColor,line width=.45pt}]
\node[noTerminal] (n0) at (0.000,0.000) {\textless{}instruccion\textgreater{}};
\node[noTerminal] (n1) at (0.000,-1.080) {\textless{}asignacion\_aritmetica\textgreater{}};
\node[terminal] (n2) at (-5.500,-2.160) {ID[z]};
\node[terminal] (n3) at (-3.300,-2.160) {=};
\node[noTerminal] (n4) at (1.100,-2.160) {\textless{}expresion\textgreater{}};
\node[noTerminal] (n5) at (-1.100,-3.240) {\textless{}operando\textgreater{}};
\node[terminal] (n6) at (-1.100,-4.320) {ID[x]};
\node[noTerminal] (n7) at (1.100,-3.240) {\textless{}operador\textgreater{}};
\node[terminal] (n8) at (1.100,-4.320) {+};
\node[noTerminal] (n9) at (3.300,-3.240) {\textless{}operando\textgreater{}};
\node[terminal] (n10) at (3.300,-4.320) {ID[y]};
\node[terminal] (n11) at (5.500,-2.160) {\#};
\draw (n1.south) -- (n2.north);
\draw (n1.south) -- (n3.north);
\draw (n5.south) -- (n6.north);
\draw (n4.south) -- (n5.north);
\draw (n7.south) -- (n8.north);
\draw (n4.south) -- (n7.north);
\draw (n9.south) -- (n10.north);
\draw (n4.south) -- (n9.north);
\draw (n1.south) -- (n4.north);
\draw (n1.south) -- (n11.north);
\draw (n0.south) -- (n1.north);
\end{tikzpicture}
\par\vspace{3pt}{\fontsize{8}{10}\selectfont\color{Muted}Representación TikZ autónoma del árbol; terminales en verde y no terminales en azul.}
\end{center}

Árbol generado desde la misma estructura que devuelve el validador. El \href{https://drive.google.com/drive/folders/1m0glvzrPN8F7qQWRx1TLlZRkeD7HSk1k}{archivo DOT editable y sus representaciones SVG y PNG} están adjuntos en Google Drive y fueron creados con Graphviz. La fuente LaTeX dibuja esa misma estructura con TikZ, sin necesitar imágenes externas. No se utilizó JFLAP.
\par

\ReportHeading{13. Relación entre raíz, nodos y hojas}

\begin{itemize}[leftmargin=13pt,itemsep=3pt,topsep=3pt,parsep=0pt]

\item La raíz es \textless{}instruccion\textgreater{}, porque el árbol representa una instrucción aislada. En un árbol del programa completo la raíz sería \textless{}programa\textgreater{}.

\item Los nodos internos son no terminales. Cada conjunto de hijos corresponde exactamente al lado derecho de una producción aplicada a su padre.

\item Las hojas, de izquierda a derecha, son ID[z], =, ID[x], +, ID[y], \#. Al conservar sus lexemas se obtiene z = x + y \#.

\item Los corchetes identifican los lexemas asociados a ID; no son símbolos adicionales de entrada.

\end{itemize}

\ReportHeading{14. Comprobación de las tres cadenas}

\begingroup\fontsize{8.0}{10.8}\selectfont
\setlength{\tabcolsep}{4pt}\renewcommand{\arraystretch}{1.16}
\arrayrulecolor{RuleColor}
\noindent\begin{tabularx}{\linewidth}{|>{\hsize=0.840000\hsize\linewidth=\hsize\raggedright\arraybackslash}X|>{\hsize=0.660000\hsize\linewidth=\hsize\raggedright\arraybackslash}X|>{\hsize=1.500000\hsize\linewidth=\hsize\raggedright\arraybackslash}X|}\hline
\rowcolor{TableHead} \textbf{Cadena} & \textbf{¿Se genera como instrucción?} & \textbf{Justificación} \\\hline
x = 20 \# & Sí & \textless{}asignacion\_\allowbreak{}simple\textgreater{}, con ID[x] y \textless{}operando\textgreater{} \ensuremath{\to} NUM[20]. \\\hline
z = x * y \# & Sí & \textless{}asignacion\_\allowbreak{}aritmetica\textgreater{}, con operandos ID[x], ID[y] y \textless{}operador\textgreater{} \ensuremath{\to} *. \\\hline
salida temporal \# & Sí & \textless{}salida\textgreater{} \ensuremath{\to} salida ID[temporal] \#. \\\hline
\end{tabularx}\par\endgroup\vspace{5pt}

\textit{Respaldo: enunciado, PDF p. 3, puntos 12-14; árbol y justificaciones derivados de la gramática propuesta.}
\par

\clearpage

\ReportSection{2.4 Pregunta 4 - Validación y detección de errores}

Los casos 1-7 se comprueban como instrucciones. El caso 8 se comprueba como programa completo. La tabla indica tanto la aceptación como el elemento que falta o está fuera de lugar.
\par

\begingroup\fontsize{8.0}{10.8}\selectfont
\setlength{\tabcolsep}{4pt}\renewcommand{\arraystretch}{1.16}
\arrayrulecolor{RuleColor}
\noindent\begin{tabularx}{\linewidth}{|>{\hsize=0.280000\hsize\linewidth=\hsize\raggedright\arraybackslash}X|>{\hsize=1.120000\hsize\linewidth=\hsize\raggedright\arraybackslash}X|>{\hsize=0.720000\hsize\linewidth=\hsize\raggedright\arraybackslash}X|>{\hsize=1.880000\hsize\linewidth=\hsize\raggedright\arraybackslash}X|}\hline
\rowcolor{TableHead} \textbf{Caso} & \textbf{Cadena} & \textbf{Resultado} & \textbf{Análisis} \\\hline
1 & x = 10 \# & Válida & Asignación simple con operando NUM y terminador \#. \\\hline
2 & z = x + y \# & Válida & Asignación aritmética con dos operandos y el operador +. \\\hline
3 & salida z \# & Válida & La salida contiene un identificador y termina en \#. \\\hline
4 & x 10 \# & Inválida & Falta = entre el identificador y el valor. \\\hline
5 & z = + x y \# & Inválida & Después de = se requiere un operando; + está antes del primer operando. \\\hline
6 & salida \# & Inválida & Falta el identificador que debe seguir a salida. \\\hline
7 & x = 20 & Inválida & Falta el terminador \#. \\\hline
8 & Begin /\allowbreak{} Programa prueba /\allowbreak{} x = 5 \# /\allowbreak{} salida x \# /\allowbreak{} End & Válida con saltos de línea; inválida con /\allowbreak{} literal & Válida al interpretar las barras como separadores editoriales de líneas. Inválida si las barras son caracteres literales del programa. \\\hline
\end{tabularx}\par\endgroup\vspace{5pt}

\ReportHeading{La doble lectura del caso 8}

Se conserva la cadena original en el archivo de casos y se añade un ejemplo separado con saltos de línea. La barra /\allowbreak{} no se transforma globalmente: dentro de una expresión representa división y dentro de /\allowbreak{}* ... */\allowbreak{} participa en la delimitación del comentario.
\par

\begingroup\setlength{\fboxsep}{6pt}
\noindent\colorbox{CodeBg}{%
\begin{minipage}{\dimexpr\linewidth-2\fboxsep\relax}
\ttfamily\fontsize{8.0}{10.5}\selectfont\raggedright
\setlength{\parindent}{0pt}\setlength{\parskip}{0pt}
\strut{}Begin\par
\strut{}Programa prueba\par
\strut{}x = 5 \#\par
\strut{}salida x \#\par
\strut{}End\par
\end{minipage}}\par\endgroup\vspace{5pt}

La versión anterior con líneas es válida. La versión literal Begin /\allowbreak{} Programa ... se rechaza al encontrar /\allowbreak{} donde la producción exige Programa. Ambas comprobaciones aparecen por separado en la evidencia.
\par

\textit{Respaldo: tabla de ocho casos del enunciado, PDF p. 3; resultados en evidencias/\allowbreak{}validacion.json.}
\par

\clearpage

\ReportSection{2.5 Extensión creativa - Declaraciones de variables}

Se elige la \textbf{opción C}: añadir declaraciones sencillas. Su utilidad sintáctica es permitir que el programa exprese un tipo y un nombre antes de las asignaciones.
\par

\ReportHeading{Sintaxis y modificación formal}

\begingroup\setlength{\fboxsep}{6pt}
\noindent\colorbox{CodeBg}{%
\begin{minipage}{\dimexpr\linewidth-2\fboxsep\relax}
\ttfamily\fontsize{8.0}{10.5}\selectfont\raggedright
\setlength{\parindent}{0pt}\setlength{\parskip}{0pt}
\strut{}\textless{}instruccion\textgreater{} ::= \textless{}asignacion\_simple\textgreater{} \textbar{} \textless{}asignacion\_aritmetica\textgreater{}\par
\hspace*{8.400em}\strut{}\textbar{} \textless{}salida\textgreater{} \textbar{} \textless{}declaracion\textgreater{}\par
\strut{}\textless{}declaracion\textgreater{} ::= \textless{}tipo\textgreater{} ID "\#"\par
\strut{}\textless{}tipo\textgreater{} ::= "entero" \textbar{} "real"\par
\end{minipage}}\par\endgroup\vspace{5pt}

Los nuevos terminales son entero y real. Los nuevos no terminales son \textless{}declaracion\textgreater{} y \textless{}tipo\textgreater{}. \textless{}declaracion\textgreater{} construye la instrucción completa; \textless{}tipo\textgreater{} elige uno de los dos tipos.
\par

\begingroup\setlength{\fboxsep}{6pt}
\noindent\colorbox{CodeBg}{%
\begin{minipage}{\dimexpr\linewidth-2\fboxsep\relax}
\ttfamily\fontsize{8.0}{10.5}\selectfont\raggedright
\setlength{\parindent}{0pt}\setlength{\parskip}{0pt}
\strut{}G\_ext = (V\_ext, \ensuremath{\Sigma}\_ext, P\_ext, S)\par
\strut{}V\_ext = V \ensuremath{\cup} \{\textless{}declaracion\textgreater{}, \textless{}tipo\textgreater{}\}\par
\strut{}\ensuremath{\Sigma}\_ext = \ensuremath{\Sigma} \ensuremath{\cup} \{entero, real\}\par
\strut{}S = \textless{}programa\textgreater{}\par
\end{minipage}}\par\endgroup\vspace{5pt}

P\_\allowbreak{}ext conserva las reglas base, reemplaza la alternativa de \textless{}instruccion\textgreater{} por la versión ampliada y añade las reglas de \textless{}declaracion\textgreater{} y \textless{}tipo\textgreater{} anteriores. La lista sigue exigiendo al menos una instrucción.
\par

\ReportHeading{Cadena válida e inválida}

\begingroup\fontsize{8.0}{10.8}\selectfont
\setlength{\tabcolsep}{4pt}\renewcommand{\arraystretch}{1.16}
\arrayrulecolor{RuleColor}
\noindent\begin{tabularx}{\linewidth}{|>{\hsize=0.840000\hsize\linewidth=\hsize\raggedright\arraybackslash}X|>{\hsize=0.540000\hsize\linewidth=\hsize\raggedright\arraybackslash}X|>{\hsize=1.620000\hsize\linewidth=\hsize\raggedright\arraybackslash}X|}\hline
\rowcolor{TableHead} \textbf{Cadena} & \textbf{Resultado} & \textbf{Justificación} \\\hline
entero edad \# & Válida & Tiene tipo, identificador y terminador. \\\hline
entero \# & Inválida & Falta el identificador después del tipo. \\\hline
\end{tabularx}\par\endgroup\vspace{5pt}

\ReportHeading{Derivación por la izquierda}

\begingroup\setlength{\fboxsep}{6pt}
\noindent\colorbox{CodeBg}{%
\begin{minipage}{\dimexpr\linewidth-2\fboxsep\relax}
\ttfamily\fontsize{8.0}{10.5}\selectfont\raggedright
\setlength{\parindent}{0pt}\setlength{\parskip}{0pt}
\strut{}\textless{}instruccion\textgreater{}\par
\strut{}\ensuremath{\Rightarrow} \textless{}declaracion\textgreater{}\par
\strut{}\ensuremath{\Rightarrow} \textless{}tipo\textgreater{} ID[edad] \#\par
\strut{}\ensuremath{\Rightarrow} entero ID[edad] \#\par
\end{minipage}}\par\endgroup\vspace{5pt}

La realización textual es entero edad \#. Este fragmento puede ocupar cualquier posición de \textless{}instruccion\textgreater{} dentro del programa. El archivo programa\_\allowbreak{}extendido.txt incluye declaraciones enteras y reales, asignaciones y salidas.
\par

La extensión define sintaxis, no una tabla de símbolos: no comprueba repetición de declaraciones, declaración previa ni compatibilidad de tipos. La palabra reservada real designa el tipo; 20.5 se reconoce como NUM.
\par

\textit{Respaldo: opción C y requisitos de la extensión del enunciado, PDF p. 4.}
\par

\clearpage

\ReportSection{3. Conclusiones}

La GLC propuesta distingue programas completos de instrucciones y separa asignaciones simples, asignaciones con una operación y salidas. ID y NUM permiten reutilizar las mismas producciones con distintos nombres y valores. Las veinte instrucciones del programa de referencia quedan cubiertas y los errores de los casos 4-7 se detectan por un elemento sintáctico concreto.
\par

Las derivaciones, el árbol y el validador comparten las mismas producciones. La extensión agrega declaraciones mediante dos nuevos no terminales y dos palabras reservadas, conservando la estructura del programa. Las decisiones sobre 0, comentarios y barras editoriales se declaran para que el alcance de la solución sea verificable.
\par

\ReportSection{4. Archivos adjuntos}

Todo el contenido del laboratorio, incluyendo el informe, la fuente LaTeX, los diagramas, la gramática, el validador, los programas de ejemplo y las evidencias de las pruebas, se encuentra disponible en la siguiente carpeta de Google Drive:
\par
\textbf{Enlace:} \url{https://drive.google.com/drive/folders/1bNQcTqB4XKDPTkW4AXRuGCijnVxjX7TX}
\par

\begin{itemize}[leftmargin=13pt,itemsep=3pt,topsep=3pt,parsep=0pt]
\item \href{https://drive.google.com/drive/folders/1m0glvzrPN8F7qQWRx1TLlZRkeD7HSk1k}{Árbol sintáctico: diagramas PNG y SVG y archivo DOT editable.}
\item \href{https://drive.google.com/drive/folders/1JhszkPPDjUckrhO5nUqpM5lOfFZOcu4h}{Gramática y validador: producciones, código y programas de ejemplo.}
\item \href{https://drive.google.com/drive/folders/19fGeTyS32ESlukLO6smH5J2xwzntVZuT}{Evidencias: validación de cadenas, árbol en JSON y registro de pruebas.}
\item \href{https://drive.google.com/drive/folders/1Mk49jfYtqgAefDx3ksJxTgen6ctqOeEM}{Fuente LaTeX y respaldo: archivo editable y paquete completo en ZIP.}
\end{itemize}

\ReportHeading{Evidencias entregadas}

\begin{itemize}[leftmargin=13pt,itemsep=3pt,topsep=3pt,parsep=0pt]

\item Informe HTML, PDF y fuente LaTeX autónoma con identificación formal, producciones, derivaciones, árbol, validación, extensión y conclusiones.

\item \href{https://drive.google.com/file/d/17ewYfIYOjwTXuDbU1OAsBNICsJnZHuss/view}{gramatica.bnf} contiene las nueve producciones agrupadas de la base y la modificación de la extensión C.

\item \href{https://drive.google.com/file/d/15_pCDRy1ZVEiERgstFPbCUh0AQehxTIB/view}{validador.py} contiene el análisis léxico y sintáctico; devuelve errores con fase, línea y columna.

\item \href{https://drive.google.com/file/d/11CZZWyp0hWTD__QrImJzayKnr4Ma42wF/view}{evidencias/\allowbreak{}validacion.json} registra once comprobaciones de cadenas, incluidas ambas lecturas del caso 8 y los dos ejemplos de extensión, más los dos programas completos aceptados.

\item \href{https://drive.google.com/file/d/1FfMsvVMZ8VBfOz0pYUDTUURfInJlWc4a/view}{evidencias/\allowbreak{}arbol\_\allowbreak{}z\_\allowbreak{}x\_\allowbreak{}mas\_\allowbreak{}y.json} conserva el árbol que sirve de fuente de la figura.

\item \href{https://drive.google.com/file/d/1EoAHXeO2O1gblCPRnflsFjsVPS1Lzjz5/view}{evidencias/\allowbreak{}pruebas.txt} registra la ejecución de 20 pruebas automatizadas; el código de esas pruebas se incluye en pruebas/\allowbreak{}.

\end{itemize}

\ReportHeading{Cómo repetir la comprobación}

\begingroup\setlength{\fboxsep}{6pt}
\noindent\colorbox{CodeBg}{%
\begin{minipage}{\dimexpr\linewidth-2\fboxsep\relax}
\ttfamily\fontsize{8.0}{10.5}\selectfont\raggedright
\setlength{\parindent}{0pt}\setlength{\parskip}{0pt}
\strut{}python3 validador.py ejemplos/programa\_referencia.txt\par
\strut{}python3 validador.py --fragmento --cadena "z = x + y \#"\par
\strut{}python3 validador.py ejemplos/programa\_extendido.txt --extension\par
\strut{}python3 -m unittest discover -s pruebas -v\par
\strut{}python3 herramientas/generar\_evidencias.py\par
\end{minipage}}\par\endgroup\vspace{5pt}

\ReportHeading{Fuentes}

1. \href{https://drive.google.com/file/d/1dWAGayupZWFKaRqMTrzPXSf6t30L21k8/view}{Laboratorio\_\allowbreak{}3\_\allowbreak{}enunciado.pdf}, documento de cinco páginas proporcionado para esta actividad. Pregunta 1: p. 2; pregunta 2: pp. 2-3; pregunta 3 y ocho casos: p. 3; extensión y entregables: p. 4; coherencia exigida: p. 5.
\par

2. sesion\_\allowbreak{}3\_\allowbreak{}4.pdf, archivo del curso: definición de gramática en PDF p. 3 y definición de L(G) en PDF p. 8. Enlace: \url{https://educa.blackboard.com/ultra/courses/_40866_1/file/_4454745_1?courseId=_40866_1}
\par

3. \href{https://drive.google.com/file/d/1OJxHs1ERN404OiHG-l35q1IYpkBgYAvt/view}{Rubrica\_\allowbreak{}Laboratorio\_\allowbreak{}3.pdf}, criterios de identificación formal, construcción de la GLC, derivaciones y árbol, validación y criterios transversales.
\par

Los desarrollos concretos de esta entrega son una propuesta propia; los archivos de referencia establecen el problema y sus conceptos.
\par

\clearpage

\ReportSection{Anexo: programa de referencia validado}

Transcripción del programa que aparece entre las páginas 1 y 2 del enunciado. Se excluyen el encabezado y la numeración editorial de página; el comentario se conserva.
\par

\begingroup\setlength{\fboxsep}{6pt}
\noindent\colorbox{CodeBg}{%
\begin{minipage}{\dimexpr\linewidth-2\fboxsep\relax}
\ttfamily\fontsize{8.0}{10.5}\selectfont\raggedright
\setlength{\parindent}{0pt}\setlength{\parskip}{0pt}
\strut{}Begin\par
\strut{}Programa operativo\_01\_02\par
\strut{}\par
\strut{}x = 1.345 \#\par
\strut{}y = 20 \#\par
\strut{}\par
\strut{}/* operaciones y salidas */\par
\strut{}\par
\strut{}z = x + y \#\par
\strut{}z2 = x + y \#\par
\strut{}h = z - x \#\par
\strut{}h2 = z - x \#\par
\strut{}k = h * y \#\par
\strut{}t = k / z \#\par
\strut{}m = k \% 10 \#\par
\strut{}\par
\strut{}salida z \#\par
\strut{}salida h \#\par
\strut{}salida k \#\par
\strut{}salida t \#\par
\strut{}salida m \#\par
\strut{}\par
\strut{}temporal = 0 \#\par
\strut{}temporal = x \#\par
\strut{}x = y \#\par
\strut{}y = temporal \#\par
\strut{}\par
\strut{}salida x \#\par
\strut{}salida y \#\par
\strut{}\par
\strut{}End\par
\end{minipage}}\par\endgroup\vspace{5pt}

\end{document}
